\documentclass[
	fontsize=12pt,
	paper=a4,
	numbers=noenddot,
	listof=totoc,
	listof=entryprefix,
	listof=nochaptergap,
	bibliography=totoc,
	parskip=half,
	openany
]{scrbook}

\usepackage{cmap}


\usepackage[utf8]{inputenc}
\usepackage[T1]{fontenc}
\usepackage{lmodern}
\usepackage{ngerman}


\usepackage{scrhack}
\usepackage[table]{xcolor}
\usepackage{chngcntr}
\usepackage{enumitem}


\usepackage{amsthm}


\usepackage{lipsum}
\usepackage{unicode-math}


\usepackage{mathtools}


\usepackage[
	colorlinks=false,
	linkcolor=black,
	citecolor=black,
	filecolor=black,
	urlcolor=black,
	bookmarks=true,
	bookmarksopen=true,
	bookmarksopenlevel=3,
	bookmarksnumbered,
	plainpages=false,
	pdfpagelabels=true,
	hyperfootnotes
]{hyperref}


\renewcommand{\familydefault}{\sfdefault}
\usepackage[onehalfspacing]{setspace}
\raggedbottom

\usepackage{anyfontsize}

\addtokomafont{chapter}{\sffamily\large\bfseries}
\addtokomafont{section}{\sffamily\normalsize\bfseries}
\addtokomafont{subsection}{\sffamily\normalsize\mdseries}
\addtokomafont{caption}{\sffamily\normalsize\mdseries}

\usepackage{relsize}
\renewcommand*{\UrlFont}{\ttfamily\smaller\relax}


\usepackage[
	bindingoffset=0cm,
	inner=2.5cm,
	outer=2.5cm,
	top=3cm,
	bottom=2cm
]{geometry}

\usepackage[
	headsepline,
	plainheadsepline
]{scrlayer-scrpage}
\clearpairofpagestyles
\addtokomafont{pagehead}{\normalfont}
\ohead*{\thepage}
\ihead*{\leftmark}


\usepackage{graphicx}
\graphicspath{{pictures/}}
\usepackage{float}

\renewcommand{\thefigure}{\arabic{figure}}
\counterwithout{figure}{chapter}

\usepackage{wrapfig}


\usepackage{booktabs}
\usepackage{multirow}
\usepackage{tabularx}

\renewcommand{\thetable}{\arabic{table}}
\counterwithout{table}{chapter}


\usepackage{listings}
\usepackage{beramono}

\renewcommand{\lstlistlistingname}{List of Code Listings}
\renewcommand{\lstlistingname}{Code Listing}
\newcommand{\listoflolentryname}{\lstlistingname}

\definecolor{codegreen}{rgb}{0,0.6,0}
\definecolor{codegray}{rgb}{0.5,0.5,0.5}
\definecolor{codepurple}{rgb}{0.5,0,0.33}
\definecolor{codepurblue}{rgb}{0.16,0.0,1.0}
\definecolor{backcolour}{rgb}{0.95,0.95,0.92}

\lstdefinestyle{codestyle}{
	backgroundcolor=\color{backcolour},
	commentstyle=\color{codegreen},
	keywordstyle=\bfseries\color{codepurple},
	numberstyle=\tiny\color{codegray},
	stringstyle=\color{codepurblue},
	basicstyle=\scriptsize\ttfamily,
	breakatwhitespace=false,
	breaklines=true,
	captionpos=b,
	keepspaces=true,
	numbers=left,
	numbersep=5pt,
	showspaces=false,
	showstringspaces=false,
	showtabs=false,
	tabsize=2,
	escapeinside={(*@}{@*)}
}
\lstset{style=codestyle, numberbychapter=false}



\KOMAoptions{toc=chapterentrydotfill}
\addtokomafont{chapterentry}{\normalfont}
\setuptoc{toc}{totoc}

\DeclareTOCStyleEntry[beforeskip=0cm]{chapter}{chapter}
\DeclareTOCStyleEntry[beforeskip=0cm]{section}{section}
\DeclareTOCStyleEntry[beforeskip=0cm]{default}{subsection}

\BeforeStartingTOC[lof]{\def\autodot{:}}
\BeforeStartingTOC[lot]{\def\autodot{:}}
\BeforeStartingTOC[lol]{\def\autodot{:}}


\iffalse
\usepackage{csquotes}
\setlength\bibitemsep{.5\baselineskip}
\setcounter{biburlnumpenalty}{9000}
\setcounter{biburllcpenalty}{9000}
\setcounter{biburlucpenalty}{9000}

\fi

\usepackage[printonlyused]{acronym}


\usepackage[title,titletoc]{appendix}

\newcommand{\appendixchapter}[1]{
\cleardoublepage
\pagenumbering{arabic}
\renewcommand{\thepage}{\thechapter-\arabic{page}}
\chapter{#1}
}

\newcommand{\monthlyreport}[2]{
\section{#1}
\centering
\includegraphics[trim=55 35 55 35,clip,width=1\textwidth]{#2}
\clearpage
}



\newtheorem{beispiel}{Beispiel}

\newtheorem{these}{These}

\newtheorem{definition}{Definition}


\usepackage[nameinlink, noabbrev]{cleveref}


\lstset{
	extendedchars=true,
	inputencoding=utf8,
	literate=
		{ä}{{\"a}}1
		{ö}{{\"o}}1
		{ü}{{\"u}}1
		{ß}{{\ss}}1
		{Ä}{{\"A}}1
		{Ö}{{\"O}}1
		{Ü}{{\"U}}1
}

\begin{document}

\begin{titlepage}
	\pagestyle{empty}

	\begin{flushright}
		\begin{figure}[ht]
			\flushright
			\includegraphics[height=2cm]{../../for_latex/hfu}
		\end{figure}
	\end{flushright}

	\begin{center}
		\vspace{3cm}

		{\fontsize{22}{22} \selectfont \textbf{Blatt 3}}\\[5mm]
		{\fontsize{18}{18} \selectfont Web-Security und Intrusion Prevention}

		\vspace{12cm}

		{\fontsize{14}{14} \selectfont Luis Staudt}
	\end{center}
\end{titlepage}


\chapter*{Aufgabe 1: Die Webanwendung}

Die Flask-Anwendung \texttt{app.py} listet ohne GET-Parameter alle Dateien im
Unterordner \texttt{files/} auf und verlinkt sie über den Parameter
\texttt{file}. Ist \texttt{file} gesetzt, gibt die Anwendung die angegebene Datei
direkt mit \texttt{send\_file} zurück:

\begin{lstlisting}[language=Python]
@app.route('/')
def index():
    if not request.args.get("file"):
        return generate_list()
    else:
        return send_file(request.args.get("file"))
\end{lstlisting}

Da der übergebene Pfad ungeprüft an \texttt{send\_file} gelangt, lassen sich
beliebige Pfade einschleusen.

\chapter*{Aufgabe 2: Die Anwendung ausnutzen}

Der Parameter darf ein absoluter Pfad sein:

\begin{lstlisting}[language=bash]
curl "http://127.0.0.1:5000/?file=/etc/passwd"
-> root:x:0:0:root:/root:/bin/bash
   daemon:x:1:1:daemon:/usr/sbin:/usr/sbin/nologin
   ...
\end{lstlisting}

Ebenso lässt sich mit \texttt{../} aus dem Arbeitsverzeichnis heraus navigieren
(Path Traversal):

\begin{lstlisting}[language=bash]
curl "http://127.0.0.1:5000/?file=files/../../../../etc/passwd"
-> root:x:0:0:root:/root:/bin/bash
   ...
\end{lstlisting}

Beide Aufrufe geben den Inhalt von \texttt{/etc/passwd} preis. Analog lassen sich
andere Dateien auslesen, etwa \texttt{app.py} selbst.

\chapter*{Aufgabe 3: Absicherung}

\section*{Erster Versuch: Präfix-Prüfung}

Geprüft wird, ob der Pfad mit \texttt{files/} beginnt:

\begin{lstlisting}[language=Python]
if not file.startswith("files/"):
    abort(403)
return send_file(file)
\end{lstlisting}

Diese Prüfung blockiert absolute Pfade (\texttt{/etc/passwd} $\rightarrow$ 403),
ist aber unzureichend: Ein relativer Pfad beginnt ebenfalls mit \texttt{files/}
und umgeht die Prüfung trotzdem.

\begin{lstlisting}[language=bash]
curl -o /dev/null -w "%{http_code}" "http://.../?file=/etc/passwd"
-> 403   (blockiert)
curl "http://.../?file=files/../../../../etc/passwd"
-> 200   (Datei wird weiterhin ausgegeben)
\end{lstlisting}

Die Ursache sind relative Pfade: Das Präfix \texttt{files/} wird durch
\texttt{../} sofort wieder verlassen.

\section*{Sichere Lösung mit os.path.realpath()}

Der übergebene Pfad wird vollständig aufgelöst, wobei \texttt{..} und Symlinks
entfernt werden, und muss tatsächlich innerhalb des erlaubten Verzeichnisses
liegen:

\begin{lstlisting}[language=Python]
import os
BASE_DIR = os.path.realpath("files")

@app.route('/')
def index():
    file = request.args.get("file")
    if not file:
        return "<h1>Keine Datei angegeben</h1>"
    requested = os.path.realpath(file)
    if requested != BASE_DIR and not requested.startswith(BASE_DIR + os.sep):
        abort(403)
    return send_file(requested)
\end{lstlisting}

Da \texttt{realpath} den Pfad kanonisiert, bevor er mit dem Basisverzeichnis
verglichen wird, scheitern sowohl absolute als auch relative Traversal-Pfade.
Weil Symlinks mit aufgelöst werden, schützt die Prüfung zugleich gegen Angriffe
über symbolische Links.

\begin{lstlisting}[language=bash]
curl -o /dev/null -w "%{http_code}" "http://.../?file=/etc/passwd"
-> 403   (blockiert)
curl -o /dev/null -w "%{http_code}" "http://.../?file=files/../../../../etc/passwd"
-> 403   (blockiert)
curl "http://.../?file=files/beispiel.txt"
-> Dies ist eine harmlose Beispieldatei.   (erlaubte Datei weiterhin verfuegbar)
\end{lstlisting}

\section*{Zusammenfassung der drei Stufen}

\begin{center}
\begin{tabular}{lcc}
\toprule
 & absoluter Pfad & relativer Traversal \\
\midrule
\texttt{app.py} (verwundbar) & 200 (Leak) & 200 (Leak) \\
\texttt{app\_naiv.py} (Präfix) & 403 & 200 (Bypass) \\
\texttt{app\_secure.py} (realpath) & 403 & 403 \\
\bottomrule
\end{tabular}
\end{center}


\end{document}
